\documentclass[11pt]{article}

\usepackage{amsmath, amsthm, amssymb}
\usepackage{parskip}
\usepackage{graphicx}
\usepackage{listings}
 \usepackage{bm}

\DeclareMathOperator{\erfc}{erfc}
\DeclareMathOperator{\Q}{Q}
\newcommand{\bvec}[1]{\bm{#1}}

\title{Computing the Areas of IAU Constellations}
\author{Dale Dickerhoof}

\date{\today}
\begin{document}
\maketitle

\begin{abstract}
The areas of the 88 constellations recognized by the International Astronomical Union (IAU) are computed
through the use of Stoke's theorem.
\end{abstract}

\section{Formulation of the Problem}

The computations here are carried out in a standard spherical coordinate system defined by a radial distance $r$, a polar angle $\theta$,
and an azimuthal angle $\phi$.

The constellation boundaries are specified in coordinates of right ascension $\alpha$ and declination $\delta$. These are
essentially a kind of spherical coordinates, with an arbitrary radial distance.  Instead of a polar angle, declination is
measured from the equator, so that $\theta=\pi/2 - \delta$.  We set $\phi=\alpha$ with all angles in radians.

Since we wish to compute the area within a boundary on the sphere, we utilize Stoke's theorem to convert the surface integral
into a line integral along the constellation boundaries (\ref{eq:stokes}).

\begin{equation}
	\iint_S \left( \nabla \times \bvec{A} \right) \cdot d\bvec{S} = \oint_\ell \bvec{A} \cdot d\bvec{\ell}
	\label{eq:stokes}
\end{equation}

In order to utilize this to compute areas within arbitrary boundaries, a vector field $\bvec{A}$ needs to be found such that
$\nabla \times \bvec{A}$ has a constant radial component on the surface of the unit sphere.  Then, the line integral in
(\ref{eq:stokes}) along the closed constellation boundaries will be proportional to the surface area within.

In spherical coordinates, the radial component of the curl is expressed as

\begin{equation}
	\left( \nabla \times \bvec{A}  \right) \cdot \bvec{\hat{r}} = \frac{1}{r \sin\theta}
		\left(
			\frac{\partial}{\partial \theta} \left( A_\phi \sin\theta \right) - \frac{\partial}{\partial\phi}A_\theta
		\right).
	\label{eq:radialcurl}
\end{equation}

\section{Evaluation of Candidate Vector Fields}

Two different choices for suitable vector fields were evaluated. Only the $A_\theta$ and $A_\phi$ components need be specified,
since $A_r$ does not appear in (\ref{eq:radialcurl}).

One choice for $\bvec{A}$ is

\begin{equation}
\begin{split}
	A_\phi &= \frac{1}{\sin\theta} \\
	A_\theta &= -r \phi \sin\theta.
	\label{eq:choice1}
\end{split}
\end{equation}

It is easily shown that the radial component of the curl is unity on the unit sphere.

The constellation boundaries for the 1875 epoch run along paths of constant $\theta$ and $\phi$ values. To evaluate
the line integrals from (\ref{eq:stokes}), we also need to evaluate $d\bvec{\ell}$.  For the choice of $\bvec{A}$ in
(\ref{eq:choice1}), paths with constant values of $\phi$ have $d\ell = r d\theta$, and can be evaluated as

\begin{equation}
	\int_{\theta_1}^{\theta_2}  \bvec{A} \cdot d\bvec{\ell} = r^2 \phi \left( \cos\theta_2 - \cos\theta_1 \right),
	\label{eq:choice1theta}
\end{equation}

setting $r=1$.

Similarly, paths of constant $\theta$ have $d\ell = r \sin\theta d\phi$, and can be evaluated as

\begin{equation}
	\int_{\phi_1}^{\phi_2} \bvec{A} \cdot d\bvec{\ell} = r \left( \phi_2 - \phi_1 \right).
	\label{eq:choice1phi}
\end{equation}

Note that in (\ref{eq:choice1phi}), the line integral does not depend on $\theta$.  When integrating the closed
constellation boundaries, the sum of the constant-$\theta$ paths goes to zero, with the exception of the
constellations containing the north and south poles.

The individual line integral terms due to (\ref{eq:choice1theta}) are computing areas from the path back to the $\phi=0$
meridian.  This makes it tricky to properly account for constellations that cross over $\phi=0$.

Another choice for $\bvec{A}$ was then considered:

\begin{equation}
\begin{split}
	A_\phi &= - \frac{\cos\theta}{\sin\theta} \\
	A_\theta &= 0.
	\label{eq:choice2}
\end{split}
\end{equation}

The constant-$\theta$ paths are evaluated as

\begin{equation}
	\int_{\phi_1}^{\phi_2} A_\phi r \sin\theta d\phi = -r \cos\theta \left( \phi_2 - \phi_1 \right),
	\label{eq:choice2theta}
\end{equation}

and the constant-$\phi$ paths are conveniently always zero.

With this choice for $\bvec{A}$ in (\ref{eq:choice2}), individual line integral segments of the boundaries are computing
areas from the constant-$\theta$ paths down to the equator.  This works without major difficulties for all but two
constellations.  For the two constellations that contain the poles, the area computed
is the area outside of the constellation to the equator, which is easily recified by the addition of the area of a half-
sphere.

The main complication that occurs in the evaluation of (\ref{eq:choice2theta}) for all constellations is that if $|\phi_2 - \phi_1| > \pi$,
it means that the shortest path goes through $\phi=0$, since $0 \le \phi < 2\pi$.  To correctly evaluate the integral,
$2\pi$ should be subtracted from the larger value of $\phi_1$ and $\phi_2$.

\begin{thebibliography}{9}

\bibitem{vi49}
Constellation Boundary Data: VI-49, http://cdsarc.u-strasbg.fr/viz-bin/cat/VI/49

\end{thebibliography}
\end{document}
